\chapter{晶体结构与性质} \label{chap:晶体结构与性质}

    ——具有规则几何外形和固定熔点的固体

    \section{各类晶体}\label{sect:各类晶体}

        \subsection{离子晶体}\label{subsec:离子晶体}

            ——正负离子通过离子键结合而成的晶体。

            基本微粒：\textcolor{red}{正、负离子}

            \(\to\) 离子键

            \(\Rightarrow\) 既无方向性，也无饱和性

            \(\Rightarrow\) 离子化合物

            \textcolor{red}{影响离子晶体熔点的因素：}

            \Circled{1} 电荷\quad \Circled{2} 半径

            影响程度：电荷 \(>\) 半径

            （典型）离子晶体的性质特征：

            \begin{enumerate}
                \item 熔沸点高且难挥发
                \item 硬度较大，不具有延展性
                \item 固态时不导电，熔融状态和溶于水后能导电
                \item 大多数易溶于极性溶剂，难溶于非极性溶剂
            \end{enumerate}

            \textcolor{blue}{离子越大，含原子越多，熔点越低。}

            硝酸乙基铵（\cemh{[C2H5NH3]NO3}）熔点：\SI{12}{\degreeCelsius}

            \(\Rightarrow\) 离子液体，一般由较大的离子构成（原子多）

            （非典型离子晶体）

            \textcolor{blue}{离子液体一般由离子半径较大的正离子（或负离子）构成，其结构不对称，较难在空间紧密堆积，正、负离子之间的静电作用小，熔点低。}

            \textcolor{blue}{有些离子晶体中还存在共价键（如 \cemh{SO4^{2-}}、\cemh{NO3-} 等多原子离子），有些还存在氢键、范德华力或配位键等作用力。}

        \subsection{共价晶体}\label{subsec:共价晶体}

            ——相邻原子之间通过共价键而形成的三维空间网状结构的晶体。

            基本微粒：\textcolor{red}{原子}\qquad 共价键

            常见物质：金刚石、晶体硅（\cemh{Si}）、水晶（\cemh{SiO2}）和金刚砂（\cemh{SiC}）

            单质硼、锗（\cemh{Ge}）、氮化硼（\cemh{BN}）、氮化硅（\cemh{Si3N4}）。

            共价晶体的性质特征：

            \begin{enumerate}
                \item 熔沸点很高，硬度很大（共价键长）
                \item 一般不导电，但 \cemh{Si}、\cemh{Ge} 半导体。
                \item 难溶于一般溶剂
            \end{enumerate}

            \textcolor{red}{影响共价晶体熔沸点和硬度的因素：共价键的键能}

            结构相似的共价晶体：原子半径越小，键长越短，键能越大，越硬、熔点越高。

            \textcolor{blue}{由于共价键有方向性，共价晶体中没有密堆积的特征（与离子晶体和金属晶体不同）。}

        \subsection{分子晶体}\label{subsec:分子晶体}

            ——只含分子的晶体。

            ——分子通过分子间作用力形成的晶体。

            基本微粒：\textcolor{red}{分子}\qquad 分子间作用力

            常见物质：

            \begin{enumerate}
                \item 非金属单质，除 \cemh{C}、\cemh{Si}、\cemh{B}；
                \item 共价化合物，除属于共价晶体的；
                \item 稀有气体
            \end{enumerate}

            含有的作用力：

            分子间作用力

            原子间：共价键（除稀有气体）

            分子晶体的性质特征：

            \begin{enumerate}
                \item 熔沸点低，有很多常温下是气体或液体，液体多易挥发，固体有些易升华（干冰、\cemh{I2}、萘）。
                \item 硬度小
                \item 不导电，固态、熔融状态均不导电，少数溶于水导电（酸）
                \item 溶解性，一般符合相似相溶。
            \end{enumerate}

            \textcolor{red}{影响分子晶体熔沸点的因素：（分子间作用力大小）}

            强\(\downarrow\)弱

            \begin{enumerate}
                \item 有无分子间氢键；
                \item 组成和结构相似，相对分子质量的大小；
                \item 相对分子质量相近时，分子极性的大小。
            \end{enumerate}

        \subsection{金属晶体}\label{subsec:金属晶体}

            \textcolor{red}{金属键}——金属正离子和自由电子，

            ——本质静电力，无饱和性、方向性。

            \textcolor{blue}{金属正离子间依靠自由电子而产生的强烈相互作用，叫做金属键。}

            金属晶体的性质特征：

            \begin{enumerate}
                \item 具有金属光泽，有良好延展性、导电性、导热性
                \item 熔点、硬度相差很大。
                \item 固态、液态均导电
            \end{enumerate}

            \textcolor{red}{影响金属晶体熔沸点的因素：（金属键强弱）}

            \begin{enumerate}
                \item 电荷数；
                \item 半径；
                \item（堆积方式）
            \end{enumerate}

            \textcolor{blue}{用「电子气理论」解释金属通性：}

            \begin{itemize}
                \item \textbf{金属光泽}：自由电子吸收所有频率的光并很快放出
                \item \textbf{导电性}：自由电子在外加电场下定向运动形成电流
                \item \textbf{导热性}：自由电子运动时与金属离子碰撞传递能量
                \item \textbf{延展性}：金属键无方向性，各层滑动后仍保持相互作用
            \end{itemize}

        \subsection{混合型晶体——石墨}\label{subsec:混合型晶体石墨}

            结构微粒：\cemh{C} 原子

            结构特点：层状结构

            作用力：

            层内：碳碳单键

            层间：范德华力

            每个 \cemh{C} 原子通过 \cemh{sp^2} 与另外 3 个 \cemh{C} 原子相连 \(\Rightarrow\) 平面网状结构。

            每个 \cemh{C} 原子还有一个垂直于平面的未杂化的 2p 轨道，有一个电子 \(\Rightarrow\) 形成覆盖整个平面的离域\(\uppi\) 键。

            层与层之间易滑动 \(\Rightarrow\) 质地软，是分子晶体特征

            层内能导电 \(\Rightarrow\) 是金属晶体特征

            层内共价键 \(\Rightarrow\) 熔点高，是共价晶体特征。

            \(\therefore\) 混合型晶体。

            \Circled{1} \cemh{sp^2} 比 \cemh{sp^3} 短

            \Circled{2} 石墨有\(\uppi\) 键

            \(\Rightarrow\) 键能：石墨 \(>\) 金刚石

        \subsection{过渡型晶体}\label{subsec:过渡型晶体}

            纯粹的晶体不多，大多数是过渡晶体。

            \cemh{Na2O}、\cemh{MgO} 偏离子晶体，当作离子晶体；

            \cemh{Al2O3} 看晶型；

            \cemh{SiO2} 偏共价晶体，当作共价晶体；

            \cemh{P2O5}、\cemh{SO3}、\cemh{Cl2O7}，分子晶体。

            \textcolor{blue}{\cemh{TiCl4} 以共价键为主的过渡型晶体，偏分子晶体（熔点\SI{-25}{\degreeCelsius}）；\cemh{TiO2} 也是共价键为主的过渡型晶体，但偏共价晶体（熔点\SI{1830}{\degreeCelsius}）。}

            \textcolor{blue}{离子液体中作用力有氢键、范德华力，处在离子晶体向分子晶体的过渡之中。}

    \section{晶体特性}\label{sect:晶体特性}

        晶体——具有规则几何外形和固定熔点的固体

        非晶体——没有固定熔点，一般也没有规则几何外形的固体。

        （是否）与颗粒大小无关。

        \subsection{自范性}\label{subsec:自范性}

            ——晶体自发地呈现封闭的、规则的几何多面体外形。

            是晶体在三维空间呈现周期有序排列的宏观表现。

            \(\Leftrightarrow\) 非晶体排列相对无序，无自范性。

        \subsection{晶体的各向异性}\label{subsec:各向异性}

            ——晶体中，不同方向上具有不同的物理性质。

            \begin{itemize}
                \item 导热性\quad 云母
                \item 导电性\quad 石墨
                \item 硬度\qquad 方解石
                \item 折光率\quad 萤石
            \end{itemize}

            不同方向排列不同 \(\Rightarrow\) 各向异性。

            \textcolor{blue}{非晶体——各向同性。}

    \section{晶胞}\label{sect:晶胞}

        ——晶体中基本的重复单元。

        找出共价键基本重复单元加以分析。

        常规的晶胞都呈平行六面体。

        \(\Rightarrow\) 数量巨大的晶胞「无隙并置」。

        \textcolor{red}{晶胞选取原则——无隙并置（不是最小）}

        「无隙」——无任何间隙；

        「并置」——平行排列，取向相同。

        同一晶体内，构成晶胞形状及原子种类、个数及几何排列是完全相同的。

        \textbf{晶胞中的配位数}

        ——指距离某个原子最近且距离相等的原子个数。

        \textbf{晶胞中粒子个数的计算}

        \(\because\) 无隙并置

        \(\therefore\) 晶胞外围的原子并不为该晶胞所独有，而是共用的。

        \begin{center}
        \begin{tabular}{|c|c|c|}
            \hline
            \textbf{位置} & \textbf{贡献份额} & \textbf{说明} \\
            \hline
            顶角 & \(1/8\) & 每个顶角原子被 8 个晶胞共用 \\
            \hline
            棱上 & \(1/4\) & 每个棱上原子被 4 个晶胞共用 \\
            \hline
            面上（面心） & \(1/2\) & 每个面心原子被 2 个晶胞共用 \\
            \hline
            内部（体心） & \(1\) & 完全属于该晶胞 \\
            \hline
        \end{tabular}
        \end{center}

        \textcolor{blue}{以上适用于平行六面体晶胞。}

        \subsection{典型金属晶体晶胞}\label{subsec:典型金属晶体晶胞}

            \subsubsection{简单立方堆积（非密堆积）}

                % [IMAGE: 简单立方堆积晶胞示意图]

                \Circled{1} 含有原子数：1 = \(8 \times 1/8\) 个

                \Circled{2} 配位数：6（简单立方配位数一定是 6）

                （\Circled{3} 空间利用率：\(52.36\%\)）

                \Circled{4} 代表金属：钋（\cemh{Po}）

            \subsubsection{体心立方堆积（密堆积）}

                % [IMAGE: 体心立方堆积晶胞示意图]

                \Circled{1} 含有原子数：2 个

                \Circled{2} 配位数：8（体心立方配位数一定是 8）

                （\Circled{3} 空间利用率：\(68.02\%\)）

                \Circled{4} 代表金属：\cemh{Na}、\cemh{K}、\cemh{Fe} 等

            \subsubsection{面心立方堆积（最密堆积）}

                % [IMAGE: 面心立方堆积晶胞示意图]

                \Circled{1} 含有原子数：4 个

                \Circled{2} 配位数：12（面心立方配位数一定是 12）

                （\Circled{3} 空间利用率：\(74.05\%\)）

                \Circled{4} 代表金属：\cemh{Cu}、\cemh{Ag}、\cemh{Au} 等

            \subsubsection{六方堆积（最密堆积）}

                % [IMAGE: 六方堆积晶胞示意图]

                \Circled{1} 含有原子数：6 个

                \Circled{2} 配位数：12（六方配位数一定是 12）

                （\Circled{3} 空间利用率：\(74.05\%\)）

                \Circled{4} 代表金属：\cemh{Mg}、\cemh{Zn}、\cemh{Ti}

        \subsection{常见离子晶体晶胞}\label{subsec:常见离子晶体晶胞}

            \subsubsection{\texorpdfstring{\cemh{NaCl}}{NaCl} 晶体}

                % [IMAGE: NaCl 晶胞示意图]

                \Circled{1} 晶胞中含 \cemh{Na+} 4 个，\cemh{Cl-} 4 个；

                \(\therefore\) \cemh{NaCl} 化学式为 \cemh{NaCl}

                \Circled{2} 每个 \cemh{Na+} 周围最近的 \cemh{Cl-} 6 个；

                每个 \cemh{Cl-} 周围最近的 \cemh{Na+} 6 个。

                \textcolor{blue}{离子化合物的化学式只表示最简个数比，不写分子式。}

                \Circled{3} 每个 \cemh{Na+} 周围最近的 \cemh{Na+} 12 个；

                每个 \cemh{Cl-} 周围最近的 \cemh{Cl-} 12 个。

            \subsubsection{\texorpdfstring{\cemh{CsCl}}{CsCl} 晶体}

                % [IMAGE: CsCl 晶胞示意图]

                \Circled{1} 晶胞中含 \cemh{Cs+} 1 个，\cemh{Cl-} 1 个；

                \Circled{2} 每个 \cemh{Cs+} 周围最近的 \cemh{Cl-} 8 个；

                每个 \cemh{Cl-} 周围最近的 \cemh{Cs+} 8 个；

                \Circled{3} 每个 \cemh{Cs+} 周围最近的 \cemh{Cs+} 6 个；

                每个 \cemh{Cl-} 周围最近的 \cemh{Cl-} 6 个。

            \subsubsection{\texorpdfstring{\cemh{CaF2}}{CaF2} 晶体}

                % [IMAGE: CaF₂ 晶胞示意图]

                \Circled{1} 晶胞中含 \cemh{Ca^{2+}} 4 个，\cemh{F-} 8 个；

                \Circled{2} 每个 \cemh{Ca^{2+}} 周围最近的 \cemh{F-} 8 个；

                每个 \cemh{F-} 周围最近的 \cemh{Ca^{2+}} 4 个；

                \Circled{3} 每个 \cemh{Ca^{2+}} 周围最近的 \cemh{Ca^{2+}} 12 个；

                每个 \cemh{F-} 周围最近的 \cemh{F-} 6 个。

            \subsubsection{\texorpdfstring{\cemh{ZnS}}{ZnS} 晶体}

                % [IMAGE: ZnS 晶胞示意图]

                \Circled{1} 晶胞中含 \cemh{Zn^{2+}} 4 个，\cemh{S^{2-}} 4 个；

                \Circled{2} 每个 \cemh{Zn^{2+}} 周围最近的 \cemh{S^{2-}} 4 个；

                每个 \cemh{S^{2-}} 周围最近的 \cemh{Zn^{2+}} 4 个；

                \Circled{3} 每个 \cemh{Zn^{2+}} 周围最近的 \cemh{Zn^{2+}} 12 个；

                每个 \cemh{S^{2-}} 周围最近的 \cemh{S^{2-}} 12 个。

        \subsection{晶体密度计算}\label{subsec:晶体密度计算}

            \subsubsection{已知晶胞边长}

                \cemh{CsCl} 中 \cemh{Cs+} \(\to\) \cemh{Cs+}：\SI{4.1E-10}{\m}，计算密度（\si{\g\per\cubic\cm}）

                \[
                    \rho = \frac{m}{V}
                    = \frac{\dfrac{M(\cemh{CsCl})}{N_\text{A}} \times \text{晶胞中个数}}{a^3}
                    = \SI{4.06}{\g\per\cubic\cm}
                \]

            \subsubsection{已知微粒半径}

                \(r(\cemh{Na+}) = \SI{101}{\pm}\)，\(r(\cemh{Cl-}) = \SI{181}{\pm}\)，计算 \(\rho\)。

                \[
                    \rho = \frac{m}{V}
                    = \frac{\dfrac{M(\cemh{NaCl})}{N_\text{A}} \times \text{个数}}
                    {\left[2(r_{\cemh{Na+}} + r_{\cemh{Cl-}})\right]^3}
                    = \SI{2.17}{\g\per\cubic\cm}
                \]

            \textcolor{blue}{晶体密度计算方法小结：根据晶胞结构确定各种微粒的数目 \(\to\) 晶胞质量 \(\to\) 晶体密度 \(=\) 晶胞质量 \(\div\) 晶胞体积。}

        \subsection{常见共价晶体晶胞}\label{subsec:常见共价晶体晶胞}

            \subsubsection{金刚石结构特点}

                \Circled{1} 每个 \cemh{C} 连另外 4 个 \cemh{C}，一起形成正四面体。

                \Circled{2} 每个 \cemh{C} 都是 \cemh{sp^3} 杂化，键长 \SI{154}{\pm}，键角\ang{109;28}。

                \Circled{3} \SI{12}{\g} 金刚石中含 \cemh{C} \(N_\text{A}\) 个，\cemh{C-C} 键 \(2N_\text{A}\) 根。

                \textcolor{blue}{碳原子数目 : \cemh{C-C} 键数目 = 1 : 2}

            \subsubsection{金刚石晶胞}

                % [IMAGE: 金刚石晶胞示意图]

                \Circled{1} 晶胞中 \cemh{C}：8 个

                \Circled{2} 每个 \cemh{C} 的配位数：4。

                同 \cemh{ZnS}

            \subsubsection{金刚石模型演变}

                （1）晶体硅、单质锗

                （2）金刚砂（\cemh{SiC}）、氮化硼（\cemh{BN}）

                \cemh{BN}：

                * 金刚石结构
                * 石墨结构

                % [IMAGE: 金刚石模型演变示意图]

            \subsubsection{\texorpdfstring{\cemh{SiO2}}{SiO2} 晶体结构特点}

                \Circled{1} 晶胞中含 8 个 \cemh{Si}，16 个 \cemh{O}

                \Circled{2} 每个 \cemh{Si} 连 4 个 \cemh{O}，每个 \cemh{O} 连 2 个 \cemh{Si}。

                \textcolor{blue}{（化合物形成的）共价晶体的化学式仅表示晶体中原子的最简个数比，也不是分子式。}

                \Circled{3} \SI{1}{\mole} \cemh{SiO2} 中含 \cemh{Si-O} 键 \(4N_\text{A}\) 根。

                \textcolor{red}{共价键有方向性/饱和性}

                \(\Rightarrow\) 无密堆积。

        \subsection{分子晶体晶胞}\label{subsec:分子晶体晶胞}

            \subsubsection{大多数分子晶体的晶胞特点}

                面心立方

                粒子数同：4 个分子

                以一个分子为中心，周围最多有 12 个紧邻的分子

                \(\Rightarrow\) 分子密堆积（范德华力无方向性）

                分子晶体中存在分子空间排列方向（即取向）问题。

                干冰：

                \(a = \SI{0.57}{\nm}\)

                \(\rho = \SI{1.578}{\g\per\cubic\cm}\)

                % [IMAGE: 干冰晶胞示意图]

            \subsubsection{冰晶体（非密堆积）}

                每个水分子依靠氢键与周围近似正四面体顶角的四个水分子连接。

                \(\Rightarrow\) 冰晶体

                \textcolor{blue}{由于氢键的方向性，使冰晶体中每个水分子与正四面体顶角的 4 个相邻水分子相互吸引，形成空隙较大的网状体。冰融化后，分子间的空隙减小，所以冰的密度比水小。}

                \textcolor{blue}{\SI{1}{\mole} \cemh{H2O} 形成的冰晶体中含氢键 \SI{2}{\mole}。}

                % [IMAGE: 冰晶胞示意图]

            \subsubsection{其他非密堆积的分子晶体}

                甲醇、乙醇

                共同点：有氢键

        \subsection{晶体 X 射线衍射分析}\label{subsec:晶体X射线衍射分析}

            一定波长 \(\Rightarrow\) 晶体会显著地产生有规律的衍射图像

            可测定：

            \begin{itemize}
                \item 晶胞的大小和形状
                \item 晶胞中原子种类、数目及在位置排列
                \item 分子中化学键的键长、键角和空间结构
            \end{itemize}

            具有不损伤样品、测量精度高、信息量大等优点。
